\section{Linear systems}

\begin{definition}[Linear-system functor]
\label{akcp-4-1-linear-system}
\dcref{4.1}
\source{MR0555258-compactifying-picard:page-0029}
For locally finitely presented $\II,\FF$ on $X/S$, define
$\LinSyst(\II,\FF)(T)$ to be the quotients $\FF_T\twoheadrightarrow\GG$
whose pseudo-ideal is isomorphic to $\II_T\otimes\LL$ for an invertible
$\LL$ on $T$.
\end{definition}

\begin{theorem}[Representability of linear systems]
\label{akcp-4-2-linear-representability}
\dcref{4.2}
\source{MR0555258-compactifying-picard:page-0029}
\uses{akcp-1-1-H,akcp-1-3-local-free,akcp-4-1-linear-system}
Let $f:X\to S$ be proper finitely presented, let $\II,\FF$ be locally finitely
presented, with $\FF$ $S$-flat, and assume that for every $T$ on which $\II_T$
is flat the canonical map
$\OO_T\to(f_T)_*\mathcal I\!som(\II_T,\II_T)$ is an isomorphism.  Then
$\LinSyst(\II,\FF)$ is represented by an open retrocompact subscheme
$U\subset\PP(H(\II,\FF))$.  Its universal quotient fits into
\[
0\to\II_U\otimes\OO_U(-1)\to\FF_U\to\GG\to0.
\]
Moreover $U=\PP(H(\II,\FF))$ iff every nonzero map
$\II(s)\to\FF(s)$ is injective on every geometric fiber.
\end{theorem}
\begin{proof}
The defining isomorphism for $H$ identifies a point of the projective bundle,
represented by a surjection $H_T\twoheadrightarrow\LL$, with a map
$\II_T\otimes\LL^{-1}\to\FF_T$ nonzero on every fiber.  A linear-system
quotient gives the same data by its chosen identification of the pseudo-ideal;
full faithfulness of $N\mapsto\II\otimes N$ makes the class independent of
choices.  Fiberwise injectivity is equivalent to flatness of the cokernel, so
the desired locus is exactly the open locus where the universal map has flat
cokernel.  The flattening theorem shows it is retrocompact, and the cokernel
gives the displayed universal sequence.  The final criterion follows by
comparing the nonzero and injective loci fiberwise.
\end{proof}

\begin{lemma}[Descent of line-bundle twists]
\label{akcp-4-3-twist-descent}
\dcref{4.3}
\source{MR0555258-compactifying-picard:page-0030}
Let $f:X\to S$ be quasi-compact and quasi-separated, $\II,\FF$ quasi-coherent,
with $\II$ locally finitely presented, and suppose
$\OO_S\simeq f_*\mathcal H\!om(\II,\II)$.  The following are equivalent:
 (a) $f_*\mathcal H\!om(\II,\FF)$ is invertible and the evaluation map
$\II\otimes f_*\mathcal H\!om(\II,\FF)\to\FF$ is an isomorphism;
 (b) $\II\otimes N\simeq\FF$ for an invertible $N$ on $S$;
 (c) $\II$ and $\FF$ are locally isomorphic over $S$;
 (d) they become isomorphic after a faithfully flat base change.
\end{lemma}
\begin{proof}
Only (d)$\Rightarrow$(a) is nonformal.  Flat base change preserves the Hom
sheaves and their adjunction maps, so after the faithfully flat cover the module
$N=f_*\mathcal H\!om(\II,\FF)$ is trivial and evaluation is an isomorphism.
Invertibility and isomorphy descend faithfully flatly.  The other implications
are immediate, and (a)$\Leftrightarrow$(b) uses the canonical evaluation map.
\end{proof}

\begin{remark}[Separatedness]
\label{akcp-4-4-separatedness}
\dcref{4.4}
\source{MR0555258-compactifying-picard:page-0031}
\uses{akcp-4-3-twist-descent}
The linear-system functor is separated for faithfully flat descent whenever the
canonical endomorphism map is an isomorphism; in particular this holds under
the hypotheses of \cref{akcp-4-2-linear-representability}.
\end{remark}

\begin{lemma}[Scalar endomorphisms]
\label{akcp-4-5-endomorphism-map}
\dcref{4.5}
\source{MR0555258-compactifying-picard:page-0032}
For a morphism of ringed spaces and an $\OO_X$-module $\II$, let
$\OO_S\to f_*\mathcal H\!om(\II,\II)$ and
$\OO_S^*\to f_*\mathcal I\!som(\II,\II)$ be the canonical maps.  The latter
is an isomorphism iff $f_*\mathcal H\!om(\II,\II)$ is invertible.  If $S$ is
locally ringed, injectivity (resp. surjectivity with locally finite generation)
of the map on the endomorphism stalk at $s$ implies injectivity (resp.
surjectivity) of the units map at $s$.
\end{lemma}
\begin{proof}
If the Hom sheaf is invertible, choose a local generator $\alpha$; the identity
is $a\alpha$ and commutativity forces $a$ a unit, so both canonical maps are
isomorphisms.  For injectivity, an element in the kernel has $1+a$ a unit and
the stalk map sends it to the identity, forcing $a=0$.  For surjectivity, the
subalgebra generated by a local section has Artinian semisimple residue algebra;
surjectivity on the residue field makes it local.  Every element is then either
a unit or differs from a unit by an element of the maximal ideal, and both lie
in the image.
\end{proof}
